\documentclass[11pt,letterpaper]{article}
\usepackage[technical-reference]{cornell-notes}
\usepackage{needspace}

\title{Clause 15: Preprocessing Directives - Cornell Notes}
\author{}
\date{}
\setDocTitle{Clause 15: Preprocessing Directives}
\setDocSubtitle{C++ 2024 Cornell Notes}
\setDocOwner{C++ 2024 Cornell Notes}
\setDocAuthor{}
\setDocDate{}
\setCornellCollection{C++ 2024 Cornell Notes}
\setCornellUnitType{Clause}
\setCornellUnitNumber{15}
\setCornellUnitTitle{Preprocessing Directives}
\setCornellCompanionLabel{C++ Technical Reference}
\hypersetup{pdftitle={Clause 15: Preprocessing Directives - Cornell Notes},pdfsubject={C++ technical study notes},pdfauthor={}}

\begin{document}
\maketitle
\begin{CornellReferenceOrientationBox}
\textbf{Purpose.} Defines conditional inclusion, source inclusion, module-related directives, macro replacement, stringizing, token pasting, rescanning, line control, diagnostics, pragmas, and predefined macros.
\par\textbf{Role.} Controls source transformation before the core syntactic grammar is applied.
\par\textbf{Principal dependencies.} Uses lexical preprocessing tokens and interacts with headers, modules, and translation phases.
\par\textbf{Primary audience.} programmers, preprocessors, build tools, and source generators.
\par\textbf{Major subject groups.} Preamble; Conditional inclusion; Source file inclusion; Module directive; Header unit importation; Macro replacement; Line control; Diagnostic directives; Pragma directive; Null directive; Predefined macro names; Pragma operator.
\end{CornellReferenceOrientationBox}
\section{Learning Objectives}
\begin{itemize}[label=\textcolor{CornellReferenceTeal}{$\blacktriangleright$}]
\item Define the governing ideas of Preamble and state the consequences for a program or implementation.
\item Identify the governing ideas of Conditional inclusion and state the consequences for a program or implementation.
\item Distinguish the governing ideas of Source file inclusion and state the consequences for a program or implementation.
\item Explain the governing ideas of Module directive and state the consequences for a program or implementation.
\item Trace the governing ideas of Header unit importation and state the consequences for a program or implementation.
\item Apply the governing ideas of Macro replacement and state the consequences for a program or implementation.
\item Analyze the governing ideas of Line control and state the consequences for a program or implementation.
\item Evaluate the governing ideas of Diagnostic directives and state the consequences for a program or implementation.
\item Diagnose the governing ideas of Pragma directive and state the consequences for a program or implementation.
\item Compare the governing ideas of Null directive and state the consequences for a program or implementation.
\end{itemize}
\section{Normative Structure Map}
\small\begin{longtable}{@{}>{\RaggedRight\arraybackslash}p{0.13\textwidth}>{\RaggedRight\arraybackslash}p{0.27\textwidth}>{\RaggedRight\arraybackslash}p{0.19\textwidth}>{\RaggedRight\arraybackslash}p{0.31\textwidth}@{}}
\toprule\textbf{Subclause} & \textbf{Subject} & \textbf{Rule category} & \textbf{Key dependencies / entities}\\\midrule\endfirsthead
\toprule\textbf{Subclause} & \textbf{Subject} & \textbf{Rule category} & \textbf{Key dependencies / entities}\\\midrule\endhead
\bottomrule\endfoot
15.1 & Preamble & core-language semantics & Uses lexical preprocessing tokens and interacts with headers, modules, and translation phases.\\
15.2 & Conditional inclusion & concurrency contract & Uses lexical preprocessing tokens and interacts with headers, modules, and translation phases.\\
15.3 & Source file inclusion & core-language semantics & Uses lexical preprocessing tokens and interacts with headers, modules, and translation phases.\\
15.4 & Module directive & core-language semantics & Uses lexical preprocessing tokens and interacts with headers, modules, and translation phases.\\
15.5 & Header unit importation & interface / declarations & Uses lexical preprocessing tokens and interacts with headers, modules, and translation phases.\\
15.6 & Macro replacement & core-language semantics & General, Argument substitution, The \# operator, The \#\# operator, and related subclauses\\
15.7 & Line control & core-language semantics & Uses lexical preprocessing tokens and interacts with headers, modules, and translation phases.\\
15.8 & Diagnostic directives & error behavior & Uses lexical preprocessing tokens and interacts with headers, modules, and translation phases.\\
15.9 & Pragma directive & core-language semantics & Uses lexical preprocessing tokens and interacts with headers, modules, and translation phases.\\
15.10 & Null directive & core-language semantics & Uses lexical preprocessing tokens and interacts with headers, modules, and translation phases.\\
15.11 & Predefined macro names & core-language semantics & Uses lexical preprocessing tokens and interacts with headers, modules, and translation phases.\\
15.12 & Pragma operator & core-language semantics & Uses lexical preprocessing tokens and interacts with headers, modules, and translation phases.\\
\end{longtable}\normalsize
\begin{CornellReferenceNormativeBox}A heading maps the clause; it does not replace the detailed conditions. For each operation, read requirements, effects, result, exceptions, complexity, remarks, and notes with their stated normative force.\end{CornellReferenceNormativeBox}
\subsection*{Rule Classification Guide}
\small\begin{longtable}{@{}>{\RaggedRight\arraybackslash}p{0.25\textwidth}>{\RaggedRight\arraybackslash}p{0.67\textwidth}@{}}\toprule\textbf{Label or category} & \textbf{How to read it}\\\midrule\endfirsthead\toprule\textbf{Label or category} & \textbf{How to read it (continued)}\\\midrule\endhead\bottomrule\endfoot
Well-formed / ill-formed & Formation is governed by syntax and semantic rules; an ill-formed program does not become well-formed because one implementation accepts it.\\
Diagnosable rule & A violation requires at least one diagnostic unless the wording places the case outside the diagnosable set.\\
No diagnostic required & The program can be ill-formed while the implementation has no diagnostic obligation.\\
Undefined behavior & No execution requirements are imposed; translation success does not establish safety or portability.\\
Unspecified behavior & One of the permitted outcomes may be chosen without documentation.\\
Implementation-defined behavior & The implementation chooses and documents the behavior.\\
Conditionally supported & The implementation may omit support but documents unsupported constructs.\\
\end{longtable}\normalsize
\section{Cornell Cue-and-Notes}
\Needspace{8\baselineskip}
\subsection*{15.1 Preamble}
\begin{longtable}{@{}>{\RaggedRight\arraybackslash\columncolor{CornellReferenceCueGray}}p{0.28\textwidth}>{\RaggedRight\arraybackslash}p{0.64\textwidth}@{}}
\rowcolor{CornellReferenceNavy}\color{white}\textbf{Cue / recall prompt} & \color{white}\textbf{Notes}\\\endfirsthead
\rowcolor{CornellReferenceNavy}\color{white}\textbf{Cue / recall prompt} & \color{white}\textbf{Notes (continued)}\\\endhead
\arrayrulecolor{CornellReferenceLineGray}
\textbf{What rule applies to Preamble?}\\[-1mm]{\scriptsize\CornellClauseRef{15.1}} & Frames the terminology, applicability, and relationships needed to interpret Preamble; detailed rules remain in the following subclauses.\\\hline
\end{longtable}
\begin{CornellReferenceSummaryBox}Preamble supplies the core-language semantics portion of this clause. The key review move is to connect its local wording to the clause-wide model, then classify each condition by normative force before relying on it.\end{CornellReferenceSummaryBox}
\Needspace{8\baselineskip}
\subsection*{15.2 Conditional inclusion}
\begin{longtable}{@{}>{\RaggedRight\arraybackslash\columncolor{CornellReferenceCueGray}}p{0.28\textwidth}>{\RaggedRight\arraybackslash}p{0.64\textwidth}@{}}
\rowcolor{CornellReferenceNavy}\color{white}\textbf{Cue / recall prompt} & \color{white}\textbf{Notes}\\\endfirsthead
\rowcolor{CornellReferenceNavy}\color{white}\textbf{Cue / recall prompt} & \color{white}\textbf{Notes (continued)}\\\endhead
\arrayrulecolor{CornellReferenceLineGray}
\textbf{What rule applies to Conditional inclusion?}\\[-1mm]{\scriptsize\CornellClauseRef{15.2}} & Defines the core-language rules for Conditional inclusion, including formation, semantic effect, interactions with type and lookup, and any ill-formed or implementation-dependent cases.\\\hline
\end{longtable}
\begin{CornellReferenceSummaryBox}Conditional inclusion supplies the concurrency contract portion of this clause. The key review move is to connect its local wording to the clause-wide model, then classify each condition by normative force before relying on it.\end{CornellReferenceSummaryBox}
\Needspace{8\baselineskip}
\subsection*{15.3 Source file inclusion}
\begin{longtable}{@{}>{\RaggedRight\arraybackslash\columncolor{CornellReferenceCueGray}}p{0.28\textwidth}>{\RaggedRight\arraybackslash}p{0.64\textwidth}@{}}
\rowcolor{CornellReferenceNavy}\color{white}\textbf{Cue / recall prompt} & \color{white}\textbf{Notes}\\\endfirsthead
\rowcolor{CornellReferenceNavy}\color{white}\textbf{Cue / recall prompt} & \color{white}\textbf{Notes (continued)}\\\endhead
\arrayrulecolor{CornellReferenceLineGray}
\textbf{What rule applies to Source file inclusion?}\\[-1mm]{\scriptsize\CornellClauseRef{15.3}} & Defines the core-language rules for Source file inclusion, including formation, semantic effect, interactions with type and lookup, and any ill-formed or implementation-dependent cases.\\\hline
\end{longtable}
\begin{CornellReferenceSummaryBox}Source file inclusion supplies the core-language semantics portion of this clause. The key review move is to connect its local wording to the clause-wide model, then classify each condition by normative force before relying on it.\end{CornellReferenceSummaryBox}
\Needspace{8\baselineskip}
\subsection*{15.4 Module directive}
\begin{longtable}{@{}>{\RaggedRight\arraybackslash\columncolor{CornellReferenceCueGray}}p{0.28\textwidth}>{\RaggedRight\arraybackslash}p{0.64\textwidth}@{}}
\rowcolor{CornellReferenceNavy}\color{white}\textbf{Cue / recall prompt} & \color{white}\textbf{Notes}\\\endfirsthead
\rowcolor{CornellReferenceNavy}\color{white}\textbf{Cue / recall prompt} & \color{white}\textbf{Notes (continued)}\\\endhead
\arrayrulecolor{CornellReferenceLineGray}
\textbf{What rule applies to Module directive?}\\[-1mm]{\scriptsize\CornellClauseRef{15.4}} & Places Module directive within the module ownership and dependency model; visibility, reachability, linkage, and textual preprocessing remain distinct.\\\hline
\end{longtable}
\begin{CornellReferenceSummaryBox}Module directive supplies the core-language semantics portion of this clause. The key review move is to connect its local wording to the clause-wide model, then classify each condition by normative force before relying on it.\end{CornellReferenceSummaryBox}
\Needspace{8\baselineskip}
\subsection*{15.5 Header unit importation}
\begin{longtable}{@{}>{\RaggedRight\arraybackslash\columncolor{CornellReferenceCueGray}}p{0.28\textwidth}>{\RaggedRight\arraybackslash}p{0.64\textwidth}@{}}
\rowcolor{CornellReferenceNavy}\color{white}\textbf{Cue / recall prompt} & \color{white}\textbf{Notes}\\\endfirsthead
\rowcolor{CornellReferenceNavy}\color{white}\textbf{Cue / recall prompt} & \color{white}\textbf{Notes (continued)}\\\endhead
\arrayrulecolor{CornellReferenceLineGray}
\textbf{What rule applies to Header unit importation?}\\[-1mm]{\scriptsize\CornellClauseRef{15.5}} & Defines the core-language rules for Header unit importation, including formation, semantic effect, interactions with type and lookup, and any ill-formed or implementation-dependent cases.\\\hline
\end{longtable}
\begin{CornellReferenceSummaryBox}Header unit importation supplies the interface / declarations portion of this clause. The key review move is to connect its local wording to the clause-wide model, then classify each condition by normative force before relying on it.\end{CornellReferenceSummaryBox}
\Needspace{8\baselineskip}
\subsection*{15.6 Macro replacement}
\begin{longtable}{@{}>{\RaggedRight\arraybackslash\columncolor{CornellReferenceCueGray}}p{0.28\textwidth}>{\RaggedRight\arraybackslash}p{0.64\textwidth}@{}}
\rowcolor{CornellReferenceNavy}\color{white}\textbf{Cue / recall prompt} & \color{white}\textbf{Notes}\\\endfirsthead
\rowcolor{CornellReferenceNavy}\color{white}\textbf{Cue / recall prompt} & \color{white}\textbf{Notes (continued)}\\\endhead
\arrayrulecolor{CornellReferenceLineGray}
\textbf{What rule applies to General?}\\[-1mm]{\scriptsize\CornellClauseRef{15.6.1}} & Frames the terminology, applicability, and relationships needed to interpret General; detailed rules remain in the following subclauses.\\\hline
\textbf{What rule applies to Argument substitution?}\\[-1mm]{\scriptsize\CornellClauseRef{15.6.2}} & Defines the core-language rules for Argument substitution, including formation, semantic effect, interactions with type and lookup, and any ill-formed or implementation-dependent cases.\\\hline
\textbf{What rule applies to The \# operator?}\\[-1mm]{\scriptsize\CornellClauseRef{15.6.3}} & Defines the core-language rules for The \# operator, including formation, semantic effect, interactions with type and lookup, and any ill-formed or implementation-dependent cases.\\\hline
\textbf{What rule applies to The \#\# operator?}\\[-1mm]{\scriptsize\CornellClauseRef{15.6.4}} & Defines the core-language rules for The \#\# operator, including formation, semantic effect, interactions with type and lookup, and any ill-formed or implementation-dependent cases.\\\hline
\textbf{What rule applies to Rescanning and further replacement?}\\[-1mm]{\scriptsize\CornellClauseRef{15.6.5}} & Defines the core-language rules for Rescanning and further replacement, including formation, semantic effect, interactions with type and lookup, and any ill-formed or implementation-dependent cases.\\\hline
\textbf{What rule applies to Scope of macro definitions?}\\[-1mm]{\scriptsize\CornellClauseRef{15.6.6}} & Locates names and entity identity for Scope of macro definitions; distinguish where a name is visible from whether declarations denote the same entity across contexts.\\\hline
\end{longtable}
\begin{CornellReferenceSummaryBox}Macro replacement supplies the core-language semantics portion of this clause. The key review move is to connect its local wording to the clause-wide model, then classify each condition by normative force before relying on it.\end{CornellReferenceSummaryBox}
\Needspace{8\baselineskip}
\subsection*{15.7 Line control}
\begin{longtable}{@{}>{\RaggedRight\arraybackslash\columncolor{CornellReferenceCueGray}}p{0.28\textwidth}>{\RaggedRight\arraybackslash}p{0.64\textwidth}@{}}
\rowcolor{CornellReferenceNavy}\color{white}\textbf{Cue / recall prompt} & \color{white}\textbf{Notes}\\\endfirsthead
\rowcolor{CornellReferenceNavy}\color{white}\textbf{Cue / recall prompt} & \color{white}\textbf{Notes (continued)}\\\endhead
\arrayrulecolor{CornellReferenceLineGray}
\textbf{What rule applies to Line control?}\\[-1mm]{\scriptsize\CornellClauseRef{15.7}} & Defines the core-language rules for Line control, including formation, semantic effect, interactions with type and lookup, and any ill-formed or implementation-dependent cases.\\\hline
\end{longtable}
\begin{CornellReferenceSummaryBox}Line control supplies the core-language semantics portion of this clause. The key review move is to connect its local wording to the clause-wide model, then classify each condition by normative force before relying on it.\end{CornellReferenceSummaryBox}
\Needspace{8\baselineskip}
\subsection*{15.8 Diagnostic directives}
\begin{longtable}{@{}>{\RaggedRight\arraybackslash\columncolor{CornellReferenceCueGray}}p{0.28\textwidth}>{\RaggedRight\arraybackslash}p{0.64\textwidth}@{}}
\rowcolor{CornellReferenceNavy}\color{white}\textbf{Cue / recall prompt} & \color{white}\textbf{Notes}\\\endfirsthead
\rowcolor{CornellReferenceNavy}\color{white}\textbf{Cue / recall prompt} & \color{white}\textbf{Notes (continued)}\\\endhead
\arrayrulecolor{CornellReferenceLineGray}
\textbf{What rule applies to Diagnostic directives?}\\[-1mm]{\scriptsize\CornellClauseRef{15.8}} & Defines the failure model for Diagnostic directives, including how an error is represented, reported, propagated, or converted into a diagnostic consequence.\\\hline
\end{longtable}
\begin{CornellReferenceSummaryBox}Diagnostic directives supplies the error behavior portion of this clause. The key review move is to connect its local wording to the clause-wide model, then classify each condition by normative force before relying on it.\end{CornellReferenceSummaryBox}
\Needspace{8\baselineskip}
\subsection*{15.9 Pragma directive}
\begin{longtable}{@{}>{\RaggedRight\arraybackslash\columncolor{CornellReferenceCueGray}}p{0.28\textwidth}>{\RaggedRight\arraybackslash}p{0.64\textwidth}@{}}
\rowcolor{CornellReferenceNavy}\color{white}\textbf{Cue / recall prompt} & \color{white}\textbf{Notes}\\\endfirsthead
\rowcolor{CornellReferenceNavy}\color{white}\textbf{Cue / recall prompt} & \color{white}\textbf{Notes (continued)}\\\endhead
\arrayrulecolor{CornellReferenceLineGray}
\textbf{What rule applies to Pragma directive?}\\[-1mm]{\scriptsize\CornellClauseRef{15.9}} & Defines the core-language rules for Pragma directive, including formation, semantic effect, interactions with type and lookup, and any ill-formed or implementation-dependent cases.\\\hline
\end{longtable}
\begin{CornellReferenceSummaryBox}Pragma directive supplies the core-language semantics portion of this clause. The key review move is to connect its local wording to the clause-wide model, then classify each condition by normative force before relying on it.\end{CornellReferenceSummaryBox}
\Needspace{8\baselineskip}
\subsection*{15.10 Null directive}
\begin{longtable}{@{}>{\RaggedRight\arraybackslash\columncolor{CornellReferenceCueGray}}p{0.28\textwidth}>{\RaggedRight\arraybackslash}p{0.64\textwidth}@{}}
\rowcolor{CornellReferenceNavy}\color{white}\textbf{Cue / recall prompt} & \color{white}\textbf{Notes}\\\endfirsthead
\rowcolor{CornellReferenceNavy}\color{white}\textbf{Cue / recall prompt} & \color{white}\textbf{Notes (continued)}\\\endhead
\arrayrulecolor{CornellReferenceLineGray}
\textbf{What rule applies to Null directive?}\\[-1mm]{\scriptsize\CornellClauseRef{15.10}} & Defines the core-language rules for Null directive, including formation, semantic effect, interactions with type and lookup, and any ill-formed or implementation-dependent cases.\\\hline
\end{longtable}
\begin{CornellReferenceSummaryBox}Null directive supplies the core-language semantics portion of this clause. The key review move is to connect its local wording to the clause-wide model, then classify each condition by normative force before relying on it.\end{CornellReferenceSummaryBox}
\Needspace{8\baselineskip}
\subsection*{15.11 Predefined macro names}
\begin{longtable}{@{}>{\RaggedRight\arraybackslash\columncolor{CornellReferenceCueGray}}p{0.28\textwidth}>{\RaggedRight\arraybackslash}p{0.64\textwidth}@{}}
\rowcolor{CornellReferenceNavy}\color{white}\textbf{Cue / recall prompt} & \color{white}\textbf{Notes}\\\endfirsthead
\rowcolor{CornellReferenceNavy}\color{white}\textbf{Cue / recall prompt} & \color{white}\textbf{Notes (continued)}\\\endhead
\arrayrulecolor{CornellReferenceLineGray}
\textbf{What rule applies to Predefined macro names?}\\[-1mm]{\scriptsize\CornellClauseRef{15.11}} & Operates on preprocessing tokens for Predefined macro names; expansion, argument substitution, rescanning, and directive context occur before core-language parsing.\\\hline
\end{longtable}
\begin{CornellReferenceSummaryBox}Predefined macro names supplies the core-language semantics portion of this clause. The key review move is to connect its local wording to the clause-wide model, then classify each condition by normative force before relying on it.\end{CornellReferenceSummaryBox}
\Needspace{8\baselineskip}
\subsection*{15.12 Pragma operator}
\begin{longtable}{@{}>{\RaggedRight\arraybackslash\columncolor{CornellReferenceCueGray}}p{0.28\textwidth}>{\RaggedRight\arraybackslash}p{0.64\textwidth}@{}}
\rowcolor{CornellReferenceNavy}\color{white}\textbf{Cue / recall prompt} & \color{white}\textbf{Notes}\\\endfirsthead
\rowcolor{CornellReferenceNavy}\color{white}\textbf{Cue / recall prompt} & \color{white}\textbf{Notes (continued)}\\\endhead
\arrayrulecolor{CornellReferenceLineGray}
\textbf{What rule applies to Pragma operator?}\\[-1mm]{\scriptsize\CornellClauseRef{15.12}} & Defines the core-language rules for Pragma operator, including formation, semantic effect, interactions with type and lookup, and any ill-formed or implementation-dependent cases.\\\hline
\end{longtable}
\begin{CornellReferenceSummaryBox}Pragma operator supplies the core-language semantics portion of this clause. The key review move is to connect its local wording to the clause-wide model, then classify each condition by normative force before relying on it.\end{CornellReferenceSummaryBox}
\section{Language-Lawyer and Library-Usage Pitfalls}
\begin{CornellReferencePitfallBox}\begin{itemize}[label=\textcolor{CornellReferenceAmber}{$\blacktriangleright$}]
\item Treating macro expansion as textual character substitution.
\item Forgetting argument prescan and rescan behavior.
\item Producing an invalid preprocessing token with token pasting.
\item Mixing module/import directives with ordinary inclusion assumptions.
\item Using reserved predefined macro names for program macros.
\item Treating a note, example, or recommended practice as though it imposed a mandatory program rule.
\end{itemize}\end{CornellReferencePitfallBox}
\section{Programmer and Implementation Checklist}
\begin{CornellReferenceChecklistBox}\begin{itemize}[label=$\square$]
\item Tokenize directive arguments before modeling expansion.
\item Track macro disablement, substitution, stringizing, pasting, and rescanning stages.
\item Evaluate conditional groups using preprocessing-expression rules.
\item Audit include and header-unit search assumptions.
\item Check required diagnostics and implementation-defined pragmas.
\item For \CornellClauseRef{15.1}, verify preamble against the precise rule category and all cited dependencies.
\item For \CornellClauseRef{15.2}, verify conditional inclusion against the precise rule category and all cited dependencies.
\item For \CornellClauseRef{15.3}, verify source file inclusion against the precise rule category and all cited dependencies.
\item For \CornellClauseRef{15.4}, verify module directive against the precise rule category and all cited dependencies.
\end{itemize}\end{CornellReferenceChecklistBox}
\section{Clause Synthesis}
\begin{CornellReferenceOrientationBox}Defines conditional inclusion, source inclusion, module-related directives, macro replacement, stringizing, token pasting, rescanning, line control, diagnostics, pragmas, and predefined macros. Controls source transformation before the core syntactic grammar is applied. A sound review distinguishes syntax and participation from runtime preconditions, keeps notes and examples non-mandatory, and follows cross-clause dependencies before making a portability claim.\end{CornellReferenceOrientationBox}
\section{Self-Test Questions}
\begin{enumerate}
\item What role does Preamble play in the clause's overall model?
\item Which normative distinctions must be kept separate when applying Conditional inclusion?
\item How would you audit a program or implementation that depends on Source file inclusion?
\item Compare Module directive with Header unit importation; what different question does each answer?
\item What role does Header unit importation play in the clause's overall model?
\item Which normative distinctions must be kept separate when applying Macro replacement?
\item How would you audit a program or implementation that depends on Line control?
\item Compare Diagnostic directives with Pragma directive; what different question does each answer?
\item Scenario: a program relies on an unstated implementation behavior while using Pragma directive. What must be checked before calling the program portable?
\item Scenario: a program relies on an unstated implementation behavior while using Null directive. What must be checked before calling the program portable?
\end{enumerate}
\clearpage\section{Answer Key}
\begin{enumerate}
\item Preamble is one part of the clause's core-language semantics model. Its role is understood by connecting the local rule in 15.1 to the clause orientation and its dependent concepts.
\item Keep formation and well-formedness separate from runtime preconditions and effects. Also distinguish mandatory wording from notes, implementation choice from undefined behavior, and interface declarations from detailed contracts; see 15.2.
\item Audit Source file inclusion by locating the controlling wording in 15.3, listing inputs and type constraints, proving any preconditions, recording effects and failure behavior, and checking lifetime, invalidation, complexity, and synchronization where applicable.
\item Module directive and Header unit importation occupy different steps or facility groups. Analyze each under its own subclause, then follow the explicit dependency between them rather than transferring one set of guarantees to the other; begin at 15.4.
\item Header unit importation is one part of the clause's interface / declarations model. Its role is understood by connecting the local rule in 15.5 to the clause orientation and its dependent concepts.
\item Keep formation and well-formedness separate from runtime preconditions and effects. Also distinguish mandatory wording from notes, implementation choice from undefined behavior, and interface declarations from detailed contracts; see 15.6.
\item Audit Line control by locating the controlling wording in 15.7, listing inputs and type constraints, proving any preconditions, recording effects and failure behavior, and checking lifetime, invalidation, complexity, and synchronization where applicable.
\item Diagnostic directives and Pragma directive occupy different steps or facility groups. Analyze each under its own subclause, then follow the explicit dependency between them rather than transferring one set of guarantees to the other; begin at 15.8.
\item First identify the exact requirement in 15.9, then classify the relied-on behavior as required, unspecified, implementation-defined, conditionally supported, or undefined. Portability requires satisfying the program obligations and avoiding dependence on a choice not guaranteed in the target environments.
\item First identify the exact requirement in 15.10, then classify the relied-on behavior as required, unspecified, implementation-defined, conditionally supported, or undefined. Portability requires satisfying the program obligations and avoiding dependence on a choice not guaranteed in the target environments.
\end{enumerate}
\section{Key-Term Recap}
\begin{longtable}{@{}>{\RaggedRight\arraybackslash}p{0.28\textwidth}>{\RaggedRight\arraybackslash}p{0.64\textwidth}@{}}\toprule\textbf{Term} & \textbf{Working meaning}\\\midrule\endfirsthead\toprule\textbf{Term} & \textbf{Working meaning (continued)}\\\midrule\endhead\bottomrule\endfoot
macro replacement & Token-based substitution driven by object-like or function-like macro definitions.\\
replacement list & The preprocessing-token sequence substituted for a macro invocation.\\
stringizing & Conversion of a macro argument's token spelling into a string literal token.\\
token pasting & Formation of a preprocessing token by concatenating adjacent token spellings.\\
conditional inclusion & Directive-controlled selection of source regions.\\
predefined macro & An implementation-provided macro with specified availability or meaning.\\
\end{longtable}
\CornellTakeaway{The preprocessor is a staged token transformer whose expansion order and directive context must be modeled exactly.}
\end{document}
